\chapter{Hemicolectomy}

\section*{Definition and Overview}
A hemicolectomy is a surgical operation to remove a section of the large bowel (colon). In a right hemicolectomy, the right side of the colon is removed and the small bowel is joined to the remaining colon; in a left hemicolectomy, the left side is removed. The operation is most commonly performed to treat colon cancer, and is also used for severe inflammatory bowel disease, polyps that cannot be removed by endoscopy, bleeding, and other conditions. Modern techniques allow most hemicolectomies to be performed laparoscopically, with shorter recovery and less pain than open surgery.

\section*{Epidemiology}
Colorectal cancer is the second most commonly diagnosed cancer in Australia, affecting around 15,000 people each year, and hemicolectomy is the standard curative treatment for colon cancer on either side of the bowel. Around 70--80\% of patients with colon cancer undergo surgery with curative intent. The operation is also performed for diverticular disease complications and inflammatory bowel disease. With screening programs detecting cancers earlier, and minimally invasive techniques, outcomes after hemicolectomy continue to improve.

\section*{Aetiology and Pathophysiology}
\begin{itemize}
    \item \textbf{Colon cancer:} the most common indication; surgery removes the tumour with an adequate margin of normal bowel and its lymph nodes
    \item \textbf{Inflammatory bowel disease:} ulcerative colitis and Crohn's disease may require resection for uncontrolled disease, strictures, or complications
    \item \textbf{Large polyps:} polyps that cannot be removed endoscopically
    \item \textbf{Diverticular disease:} complicated diverticulitis with abscess, perforation, or fistula
    \item \textbf{Ischaemic and bleeding bowel:} when bowel loses its blood supply or bleeds uncontrollably
    \item \textbf{The procedure:} the affected bowel and its blood supply and lymph nodes are removed, and the remaining bowel ends are joined (anastomosis) to restore continuity
\end{itemize}

\section*{Clinical Features}
\begin{itemize}
    \item Patients commonly present with changes in bowel habit, blood in the stool, abdominal pain, or anaemia
    \item Some colon cancers are found through screening before symptoms develop
    \item Emergency surgery may be needed for bowel obstruction, perforation, or severe bleeding
    \item After surgery, most patients have an uncomplicated recovery in hospital for several days
    \item Bowel function returns gradually, with initial loose stools settling over weeks
\end{itemize}

\section*{Differential Diagnosis}
\begin{itemize}
    \item The decision to operate depends on the diagnosis, the stage of any cancer, and the patient's fitness
    \item Some early cancers and large polyps can be removed by endoscopy alone
    \item Diverticular disease and inflammatory bowel disease may be managed medically without surgery
    \item Metastatic cancer may be treated with chemotherapy, with surgery sometimes combined or deferred
    \item For obstructing or bleeding disease, surgery may be urgent rather than elective
\end{itemize}

\section*{When to Seek Medical Care}
Anyone with new bowel symptoms, blood in the stool, or a family history of bowel cancer should discuss screening with their doctor. After surgery, report fever, worsening pain, wound problems, or inability to open the bowels.

\textbf{Call 000 (or local emergency services) for:}
\begin{itemize}
    \item Sudden severe abdominal pain, which may indicate perforation or obstruction
    \item Heavy bleeding from the rectum
    \item High fever, shaking, or collapse after surgery (possible anastomotic leak or infection)
    \item Vomiting with abdominal distension and inability to pass wind or stool
    \item Signs of a heart attack or stroke after surgery
\end{itemize}

\section*{Diagnosis and Investigations}
\begin{itemize}
    \item \textbf{Colonoscopy with biopsy:} confirms the diagnosis and location of cancer or other disease
    \item \textbf{CT scan:} stages the cancer and detects spread to lymph nodes or other organs
    \item \textbf{Blood tests:} full blood count for anaemia, carcinoembryonic antigen (CEA) tumour marker, and organ function
    \item \textbf{Preoperative assessment:} fitness for surgery, including heart, lung, and general health
    \item \textbf{Postoperative pathology:} the removed specimen is examined to confirm the cancer stage and guide further treatment
    \item \textbf{Genetic testing:} for hereditary bowel cancer syndromes where indicated
\end{itemize}

\section*{Management}
\begin{itemize}
    \item \textbf{Laparoscopic surgery:} minimally invasive removal through small incisions, with faster recovery
    \item \textbf{Open surgery:} a larger incision, used for complex cases, emergency surgery, or very large tumours
    \item \textbf{Anastomosis:} the bowel is joined to restore continuity; a temporary stoma may be created if healing is at risk
    \item \textbf{Enhanced recovery:} early eating, mobilisation, and minimal opioid use speed recovery
    \item \textbf{Adjuvant chemotherapy:} offered after surgery for cancers with high-risk features
    \item \textbf{Surveillance:} regular review, colonoscopy, and imaging after curative surgery
    \item \textbf{Bowel adaptation:} most people have normal bowel function after recovery; some need dietary adjustment
\end{itemize}

\section*{Complications}
\begin{itemize}
    \item Anastomotic leak, a serious early complication causing peritonitis
    \item Bleeding, infection, and wound complications
    \item Ileus, a temporary slowing of bowel function
    \item Bowel obstruction from adhesions, which can occur months to years later
    \item Stoma complications if a temporary stoma is formed
    \item Recurrence of cancer if not completely removed
    \item Anaesthesia and cardiovascular risks, as with any major surgery
\end{itemize}

\section*{Prognosis and Outcomes}
The outlook after hemicolectomy is excellent for most patients. For colon cancer confined to the bowel, surgery is curative in the majority of cases, and survival depends mainly on the cancer stage. Recovery from the operation itself is good, with most people leaving hospital within a week and returning to normal activities within a few weeks to months. Laparoscopic surgery has reduced pain, scarring, and recovery time. Long-term bowel function is usually normal, though some people have more frequent or looser stools.

\section*{Prevention and Self-Care}
\begin{itemize}
    \item Participate in the National Bowel Cancer Screening Program from age 45--50
    \item See a doctor promptly for blood in the stool, bowel habit changes, or unexplained anaemia
    \item Maintain a diet rich in fibre, vegetables, and fruit, and limit processed and red meat
    \item Keep physically active and maintain a healthy weight
    \item Limit alcohol and avoid smoking
    \item Follow the surgical team's recovery instructions, including wound care and activity limits
    \item Attend all follow-up appointments, scans, and colonoscopies
    \item Discuss genetic testing and family screening if there is a strong family history
\end{itemize}

\section*{Sources and Further Reading}
The following reputable sources provide further information for healthcare students and patients seeking reliable, current advice about hemicolectomy.

\begin{itemize}
    \item Bowel Cancer Australia. \textit{Colon cancer surgery: hemicolectomy.} https://www.bowelcanceraustralia.org/
    \item Cancer Council. \textit{Bowel cancer treatment: surgery.} https://www.cancer.org.au/
    \item Healthdirect. \textit{Bowel surgery for cancer.} https://www.healthdirect.gov.au/
    \item NHS. \textit{Colectomy: what to expect.} https://www.nhs.uk/conditions/colectomy/
    \item Mayo Clinic. \textit{Colectomy: why it's done and recovery.} https://www.mayoclinic.org/
    \item National Bowel Cancer Screening Program. \textit{Screening information.} https://www.health.gov.au/
\end{itemize}
